\section{Remaining Tests in the Revised Plan}
\label{sec:method-plan-validation}

The remaining campaign targets a repeatable V2 baseline, power-limit diagnosis,
and a matched V3 comparison \cite{projectOctoberPlan}. AP-01--AP-08 are defined
in Section~\ref{sec:revised-plan-criteria}; current results are summarized in
Section~\ref{sec:res-revised-plan}.

\subsection{Repeat the V2 Baseline}

The reported V2 operating point near \SI{85}{\watt} must be repeated with
the hardware, rework, magnetics, resonant values, controller target, firmware version, settings and mode
identified. AC input, output setting and negotiation, load and cable, auxiliary
supplies, cooling, ambient and duration must also be recorded. Probe scaling
and zero must be checked. If the original configuration cannot be reconstructed,
the known differences must be stated.

The measurements must identify the analyzer, ranges and accuracy,
probe positions, bandwidth, averaging and timing, and thermal sensing method.
Protocol tests must identify advertised capabilities and requested/accepted
sink settings.

\subsection{Power and Temperature Measurements}

Real AC input power, external connector voltage, and load current must be
measured over the same steady interval for the calculation in
Section~\ref{sec:method-efficiency}. Secondary-bulk measurements and the
primary-current input PA2 describe upstream control rather than connector
output power. The boundary must include internally supplied auxiliaries and
measured external supply power, with any exclusions stated.

Instrument ranges, scaling, calibration, and the propagation method in
Section~\ref{sec:method-uncertainty} define the uncertainty. Unmatched
averaging windows, changing load, or omitted auxiliary power invalidate a
matched efficiency estimate.

Continuous-power tests require a specified duration, thermal-stability
criterion, and numerical stop limits. Initial condition, elapsed time,
ambient, airflow, orientation, heat sinks, and temperature histories must be
recorded; isolated electrical and thermal captures are insufficient.

\subsection{Staged Loading and Controlled Changes}

The planned sequence in Table~\ref{tab:method-revised-load-sequence} is limited
by the selected PDO/APDO and cable, connector, load, and converter ratings.
Unsupported voltage/current combinations must be omitted or reconfigured.

\begin{table}[H]
\centering
\small
\caption[Planned load sequence.]{Planned load sequence for the remaining tests.}
\label{tab:method-revised-load-sequence}
\begin{tabularx}{\textwidth}{@{}L{0.28\textwidth}Y@{}}
\toprule
Stage & Test and condition for increasing load\\
\midrule
\SI{25}{\watt}, \SI{45}{\watt}, \SI{65}{\watt} & Measure regulation, power quality, efficiency, stress, and heating at increasing load under an applicable output setting.\\
\SI{75}{\watt}, \SI{80}{\watt}, approximately \SI{85}{\watt} & Repeat the upper baseline range and document whether the plan's V2 point is reproduced.\\
\SI{90}{\watt}, \SI{95}{\watt}, \SI{100}{\watt} & Proceed in stages only after the previous point passes the defined stress, temperature, control, and interface checks.\\
Final stable point & Complete the defined-duration thermal/stability run and the full measurement set; report the highest repeatable power if the target is not reached.\\
\bottomrule
\end{tabularx}
\end{table}

Stop limits must follow component ratings, expected stress, sensing range,
and thermal assumptions, with surface and junction limits treated separately.
Limit crossings, loss of regulation/feedback, unexplained protection, or
abnormal switching require termination and a saved event trace.

Controlled retests must change one item while retaining the remaining setup.
The log must identify the part or setting, rationale, expected effect, and
before/after records. Candidates include sensing calibration, control limits,
switching timing, resonant and charge-pump values, gate behavior, conduction
paths, rectification, and cooling. Calibration changes require raw readings
and both scales to separate display correction from electrical improvement.

\subsection{Output, Transient, and Thermal Tests}

Fixed-PDO tests must record advertised and negotiated voltage/current,
regulation, load range, and attach/detach behavior. PPS tests require an
advertised APDO, requested/accepted settings, measured connector output,
transitions, and limit behavior.

Startup, shutdown, load steps/removal, no-load recovery, and mode transitions
require synchronized connector VBUS, relevant bulk/control signals, and event
markers. Protection tests must define the applied fault, trigger, response,
and recovery within the bench limits.

At the final accepted power, efficiency, PF/THD, ripple, transients, stress,
and extended thermal stability must be assessed. Drift, protection,
interruptions, and cooldown/restart behavior belong in the run record.
A defined-duration result applies to those conditions rather than lifetime
reliability.

\subsection{Comparing V2 and V3 Power Density}
\label{sec:method-density}

For both assemblies, the rectangular envelope of the complete installed
converter would be measured:
\begin{equation}
 V_{\mathrm{assembly}}=LWH,
 \qquad
 \rho_P=\frac{P_{\mathrm{out,continuous}}}{V_{\mathrm{assembly}}}.
 \label{eq:method-assembly-density}
\end{equation}
The envelope includes both PCB sides, required daughterboards/connectors,
heat sinks, and functional hardware moved off-board. Both assemblies must
use the same axes and treatment of enclosure, cooling, and supports; bench
instrumentation is excluded.

The numerator requires externally verified continuous output under stated
duration, ambient, airflow, input, output setting, and thermal conditions.
Density must be reported with efficiency, temperature, and margin. A nominal
\SI{100}{\watt} rating and a short measured peak are not comparable under AP-07.

The approximately \SI{8}{\percent} V3 reduction is a design-volume estimate
with USB-PD omitted, not a complete-supply comparison. External regulation,
PD control, connectors, and required cooling must enter the system volume.
Verified dimensions and matched continuous-power tests remain unavailable,
so no V2/V3 density result is reported. AP-08 specifies completed V2 testing
as the fallback if V3 is not ready.
